\paragraph{碰撞频率代数曲线、奇素数 shadow 与 Rényi--2 刚性}
在碰撞势的三次压力方程、核版 RH/Ramanujan 窗口、分歧晶格以及 Fibonacci 共振常数已经建立之后，还可把频率几何、素数扭曲与均匀性障碍进一步压缩为下列五条结果。

\begin{theorem}[碰撞频率的代数曲线与精确大偏差参数化]\label{thm:xi-time-part9p-collision-frequency-algebraic-ldp}
\leanverified{paper\_xi\_time\_part9p\_collision\_frequency\_algebraic\_ldp}
沿用命题 \ref{prop:real-input-40-collision-pressure} 的记号，令
\[
u=e^\theta,\qquad
P_2(\theta):=\log \rho(e^\theta),\qquad
\alpha(\theta):=P_2'(\theta).
\]
则在正实扭曲支 \(u>0\) 上成立：
\begin{enumerate}
  \item 碰撞频率 \(\alpha\) 与扭曲参数 \(u\) 满足显式三次代数曲线
  \begin{equation}\label{eq:xi-time-part9p-collision-frequency-algebraic-curve}
  \bigl(4u^3+25u^2+88u+8\bigr)\bigl(\alpha^3-\alpha^2\bigr)
  +\bigl(u^3+10u^2+24u\bigr)\alpha-2u^2=0.
  \end{equation}
  \item 以 Perron 根 \(\rho\) 为参数，可写成
  \begin{equation}\label{eq:xi-time-part9p-collision-frequency-rho-param}
  u(\rho)=\frac{\rho(\rho^2-2\rho-1)}{\rho-1},
  \qquad
  \alpha(\rho)=\frac{(\rho^2-2\rho-1)(\rho-1)}{2\rho^3-5\rho^2+4\rho+1},
  \qquad \rho\in(1+\sqrt2,\infty).
  \end{equation}
  \item 记
  \[
  I_2(a):=\sup_{\theta\in\RR}\Bigl\{a\theta-\bigl(P_2(\theta)-P_2(0)\bigr)\Bigr\},
  \]
  则在参数点 \(a=\alpha(\rho)\) 处有闭式
  \begin{equation}\label{eq:xi-time-part9p-collision-frequency-rate-param}
  I_2\bigl(\alpha(\rho)\bigr)
  =
  \alpha(\rho)\log\!\frac{\rho(\rho^2-2\rho-1)}{\rho-1}
  -\log\!\frac{\rho}{\varphi^2}.
  \end{equation}
\end{enumerate}
\end{theorem}
\begin{proof}
记
\[
f(\rho,u):=\rho^3-2\rho^2-(u+1)\rho+u.
\]
由 \(f(\rho(u),u)=0\) 对 \(u\) 作隐函数求导，得到
\[
\frac{\mathrm d\rho}{\mathrm du}
=-\frac{\partial_u f}{\partial_\rho f}
=\frac{\rho-1}{3\rho^2-4\rho-(u+1)}.
\]
又因 \(u=e^\theta\)，故
\[
\alpha(\theta)=\frac{\mathrm d}{\mathrm d\theta}\log\rho(e^\theta)
=\frac{u}{\rho}\frac{\mathrm d\rho}{\mathrm du}.
\]
由 \(f(\rho,u)=0\) 解出
\[
u(\rho)=\frac{\rho(\rho^2-2\rho-1)}{\rho-1}.
\]
对正实扭曲支 \(u>0\)，必有 \(\rho>1+\sqrt2\)，故上式分母不为零且右端为正。将 \(u(\rho)\) 代回 \(\alpha(\theta)\) 的表达式并化简，即得
\eqref{eq:xi-time-part9p-collision-frequency-rho-param} 中的 \(\alpha(\rho)\)。

为证明 \eqref{eq:xi-time-part9p-collision-frequency-algebraic-curve}，只需消去参数 \(\rho\)。把
\[
u(\rho-1)-\rho(\rho^2-2\rho-1)=0
\]
与
\[
\alpha(2\rho^3-5\rho^2+4\rho+1)-(\rho^2-2\rho-1)(\rho-1)=0
\]
视为关于 \(\rho\) 的两条代数方程，取 Sylvester 结果式即可得到
\eqref{eq:xi-time-part9p-collision-frequency-algebraic-curve}（结果式只差一个非零常数因子，不影响零点集）。

最后，由命题 \ref{prop:real-input-40-collision-all-real} 的全实谱性与有限状态 Perron 压力的一般谱理论，\(P_2\) 在实轴上实解析且严格凸；因此 Legendre--Fenchel 对偶在驻点
\(
a=P_2'(\theta)
\)
处取到。代入
\[
\theta=\log u(\rho),
\qquad
P_2(\theta)-P_2(0)=\log\rho-\log\varphi^2
\]
便得到 \eqref{eq:xi-time-part9p-collision-frequency-rate-param}。证毕。
\end{proof}

\begin{theorem}[碰撞势 Perron 支的正实有理参数化与单调双射]\label{thm:xi-time-part9p-collision-pressure-rational-monotone-branch}
\leanverified{paper\_xi\_time\_part9p\_collision\_pressure\_rational\_monotone\_branch}
沿用定理 \ref{thm:xi-time-part9p-collision-frequency-algebraic-ldp} 的记号。正实 Perron 支满足
\[
\boxed{\;
u(\rho)=\frac{\rho(\rho^2-2\rho-1)}{\rho-1}.\;}
\]
并且
\[
\boxed{\;
u>0
\iff
\rho>1+\sqrt2.\;}
\]
进一步，
\[
\boxed{\;
u'(\rho)=\frac{2\rho^3-5\rho^2+4\rho+1}{(\rho-1)^2}>0
\qquad(\rho>1+\sqrt2).\;}
\]
因此
\[
u:(1+\sqrt2,\infty)\longrightarrow (0,\infty)
\]
为严格递增双射，其逆映射正是正实扭曲参数 \(u\mapsto \rho(u)\) 的 Perron 支。
\end{theorem}
\begin{proof}
定理 \ref{thm:xi-time-part9p-collision-frequency-algebraic-ldp} 已给出
\[
u(\rho)=\frac{\rho(\rho^2-2\rho-1)}{\rho-1}.
\]
在 \(\rho>1\) 上，分母为正，故 \(u(\rho)>0\) 当且仅当
\[
\rho^2-2\rho-1>0,
\]
也即
\[
\rho>1+\sqrt2.
\]
这就得到了正实支的精确参数区间。

再对 \(u(\rho)\) 直接求导，得到
\[
u'(\rho)=\frac{2\rho^3-5\rho^2+4\rho+1}{(\rho-1)^2}.
\]
而分子可重写为
\[
2\rho^3-5\rho^2+4\rho+1
=
(\rho^2-2\rho-1)(2\rho-1)+4\rho.
\]
当 \(\rho>1+\sqrt2\) 时，上式右端各项都严格为正，故 \(u'(\rho)>0\)。

最后，
\[
\lim_{\rho\downarrow 1+\sqrt2}u(\rho)=0,
\qquad
\lim_{\rho\to\infty}u(\rho)=\infty,
\]
再结合严格单调性，即得 \(u:(1+\sqrt2,\infty)\to(0,\infty)\) 为双射。证毕。
\end{proof}

\begin{theorem}[碰撞势压力的显式导数公式与全轴严格凸性]\label{thm:xi-time-part9p-collision-pressure-derivatives-global-convexity}
沿用定理 \ref{thm:xi-time-part9p-collision-frequency-algebraic-ldp} 与
\ref{thm:xi-time-part9p-collision-pressure-rational-monotone-branch} 的记号。若
\[
\rho=\rho(e^\theta)\in(1+\sqrt2,\infty),
\]
则
\[
\boxed{\;
P_2'(\theta)=
\frac{(\rho-1)(\rho^2-2\rho-1)}
{2\rho^3-5\rho^2+4\rho+1}.\;}
\]
并且
\[
\boxed{\;
P_2''(\theta)=
\frac{\rho(\rho-1)(\rho^2-2\rho-1)\bigl(\rho^4+4\rho^3-10\rho^2+4\rho-3\bigr)}
{\bigl(2\rho^3-5\rho^2+4\rho+1\bigr)^3}.\;}
\]
特别地，对全部 \(\theta\in\RR\) 都有
\[
\boxed{\;
P_2''(\theta)>0.\;}
\]
因此 \(P_2\) 在整个实轴上严格凸。又在 \(\theta=0\) 时 \(u=1\) 且 \(\rho=\varphi^2\)，故
\[
P_2'(0)=\frac{3-\sqrt5}{10},
\qquad
P_2''(0)=\frac{6\sqrt5-5}{125},
\]
与零点处的既有 \(\sqrt5\)-闭式完全一致。
\end{theorem}
\begin{proof}
第一式正是定理 \ref{thm:xi-time-part9p-collision-frequency-algebraic-ldp} 中的
\(\alpha(\rho)\) 参数化，因为
\[
P_2'(\theta)=\alpha(\theta)=\alpha(\rho).
\]

再对该有理函数按 \(\rho\) 求导，可得
\[
\frac{\dd}{\dd\rho}P_2'(\theta(\rho))
=
\frac{\rho^4+4\rho^3-10\rho^2+4\rho-3}
{\bigl(2\rho^3-5\rho^2+4\rho+1\bigr)^2}.
\]
另一方面，由
\[
\theta=\log u(\rho)
\]
与定理 \ref{thm:xi-time-part9p-collision-pressure-rational-monotone-branch}，
\[
\frac{\dd}{\dd\theta}
=
\frac{u(\rho)}{u'(\rho)}\frac{\dd}{\dd\rho}
=
\frac{\rho(\rho-1)(\rho^2-2\rho-1)}
{2\rho^3-5\rho^2+4\rho+1}\frac{\dd}{\dd\rho}.
\]
将两式相乘，便得到 \(P_2''(\theta)\) 的闭式表达。

为证正性，只需检查各因子在 \(\rho>1+\sqrt2\) 上都严格为正。由上一条定理，
\[
\rho>0,
\qquad
\rho-1>0,
\qquad
\rho^2-2\rho-1>0,
\qquad
2\rho^3-5\rho^2+4\rho+1>0.
\]
对新出现的四次因子，可写成
\[
\rho^4+4\rho^3-10\rho^2+4\rho-3
=
(\rho^2-2\rho-1)(\rho^2+6\rho+3)+16\rho,
\]
故在 \(\rho>1+\sqrt2\) 上亦严格为正。因此 \(P_2''(\theta)>0\) 对全部 \(\theta\in\RR\) 成立，严格凸性随之得到。

最后，在 \(\theta=0\) 时 \(u=e^\theta=1\)。由
\[
u(\rho)=\frac{\rho(\rho^2-2\rho-1)}{\rho-1}
\]
与正实 Perron 支的唯一性，必有 \(\rho=\varphi^2\)。将 \(\rho=\varphi^2\) 代入上两式并化简，即得
\[
P_2'(0)=\frac{3-\sqrt5}{10},
\qquad
P_2''(0)=\frac{6\sqrt5-5}{125}.
\]
证毕。
\end{proof}

\begin{theorem}[碰撞势上尾物理支的大偏差率函数单分支统一化]\label{thm:xi-time-part9p-collision-frequency-physical-branch-monotonicity}
沿用定理 \ref{thm:xi-time-part9p-collision-frequency-algebraic-ldp} 的记号，并记
\[
\alpha_\ast:=P_2'(0)=\frac{3-\sqrt5}{10}.
\]
在经过 \(u=1\) 的物理上尾支
\[
\rho\ge \varphi^2
\]
上，
\[
\alpha(\rho)
=
\frac{(\rho^2-2\rho-1)(\rho-1)}{2\rho^3-5\rho^2+4\rho+1}
=
\frac{\rho^3-3\rho^2+\rho+1}{2\rho^3-5\rho^2+4\rho+1}
\]
严格单调递增。因而
\[
\alpha\bigl([\varphi^2,\infty)\bigr)
=
\left[\alpha_\ast,\frac12\right),
\]
并且对每个
\(
a\in[\alpha_\ast,\frac12)
\)
恰有唯一的 \(\rho\in[\varphi^2,\infty)\) 使 \(a=\alpha(\rho)\)。对应速率函数由单一实代数参数 \(\rho\) 统一给出：
\[
I_2\bigl(\alpha(\rho)\bigr)
=
\alpha(\rho)\log\!\frac{\rho(\rho^2-2\rho-1)}{\rho-1}
-\log\!\frac{\rho}{\varphi^2},
\qquad \rho\in[\varphi^2,\infty).
\]
特别地，物理上尾支不存在内部回折或分支切换。
\end{theorem}
\begin{proof}
由定理 \ref{thm:xi-time-part9p-collision-frequency-algebraic-ldp}，
\[
\alpha(\rho)=\frac{\rho^3-3\rho^2+\rho+1}{2\rho^3-5\rho^2+4\rho+1}.
\]
直接求导得
\[
\alpha'(\rho)
=
\frac{\rho^4+4\rho^3-10\rho^2+4\rho-3}{(2\rho^3-5\rho^2+4\rho+1)^2}.
\]
记
\[
q(\rho):=\rho^4+4\rho^3-10\rho^2+4\rho-3.
\]
则
\[
q(2)=13>0,
\qquad
q'(\rho)=4\rho^3+12\rho^2-20\rho+4>0
\qquad(\rho\ge 2).
\]
故 \(q(\rho)>0\) 对所有 \(\rho\ge 2\) 成立。另一方面，记
\[
d(\rho):=2\rho^3-5\rho^2+4\rho+1.
\]
则
\[
d'(\rho)=2(3\rho-2)(\rho-1)>0
\qquad(\rho\ge \varphi^2),
\]
并且
\[
d(\varphi^2)=5\varphi^2>0.
\]
因此 \(d(\rho)>0\) 在 \([\varphi^2,\infty)\) 上恒成立，分母平方严格为正，从而
\[
\alpha'(\rho)>0
\qquad(\rho\ge \varphi^2).
\]

再由 \(u(\varphi^2)=1\) 与
\(
\alpha_\ast=P_2'(0)
\)
可知
\[
\alpha(\varphi^2)=\alpha_\ast=\frac{3-\sqrt5}{10}.
\]
同时，
\[
\alpha(\rho)
=
\frac{\rho^3-3\rho^2+\rho+1}{2\rho^3-5\rho^2+4\rho+1}
\longrightarrow \frac12
\qquad(\rho\to\infty).
\]
故其像区间恰为 \([\alpha_\ast,\frac12)\)，且每个可达 \(a\) 对应唯一参数 \(\rho\)。最后，速率函数公式正是
\eqref{eq:xi-time-part9p-collision-frequency-rate-param} 在该区间上的限制。证毕。
\end{proof}

\begin{theorem}[实正扭曲分支的全局单调坐标化]\label{thm:xi-time-part9p-positive-branch-global-coordinate}
\leanverified{paper\_xi\_time\_part9p\_positive\_branch\_global\_coordinate}
记
\[
\rho_0:=1+\sqrt2.
\]
则 \eqref{eq:xi-time-part9p-collision-frequency-rho-param} 中的参数函数在区间 \((\rho_0,\infty)\) 上满足：
\begin{enumerate}
  \item \(u(\rho)\) 严格递增，并给出双射
  \[
  u:(\rho_0,\infty)\xrightarrow{\sim}(0,\infty);
  \]
  \item \(\alpha(\rho)\) 严格递增，并给出双射
  \[
  \alpha:(\rho_0,\infty)\xrightarrow{\sim}\left(0,\frac12\right).
  \]
\end{enumerate}
从而对每个实正扭曲参数 \(u>0\)，存在唯一 \(\rho>\rho_0\) 使
\[
u=\frac{\rho(\rho^2-2\rho-1)}{\rho-1},
\]
并且存在唯一碰撞频率 \(\alpha(u)\in(0,\tfrac12)\) 与之对应。亦即，\(\rho\)、\(u\) 与物理实轴上的碰撞频率在实正支上给出同一条全局一维坐标。
\end{theorem}
\begin{proof}
由 \eqref{eq:xi-time-part9p-collision-frequency-rho-param}，
\[
u'(\rho)=\frac{2\rho^3-5\rho^2+4\rho+1}{(\rho-1)^2}.
\]
记
\[
q(\rho):=2\rho^3-5\rho^2+4\rho+1.
\]
则
\[
q'(\rho)=6\rho^2-10\rho+4=2(3\rho-2)(\rho-1)>0
\qquad(\rho>1),
\]
故 \(q\) 在 \((1,\infty)\) 上严格递增。又
\[
q(\rho_0)=4+4\sqrt2>0,
\qquad
u(\rho_0)=0,
\qquad
u(\rho)\sim \rho^2\to+\infty
\quad(\rho\to+\infty),
\]
从而对一切 \(\rho>\rho_0\) 都有 \(u'(\rho)>0\)，并且 \(u\) 把 \((\rho_0,\infty)\) 严格单调地双射到 \((0,\infty)\)。

另一方面，
\[
\alpha'(\rho)=
\frac{\rho^4+4\rho^3-10\rho^2+4\rho-3}
{(2\rho^3-5\rho^2+4\rho+1)^2}.
\]
记
\[
p(\rho):=\rho^4+4\rho^3-10\rho^2+4\rho-3.
\]
则
\[
p''(\rho)=12\rho^2+24\rho-20>0
\qquad(\rho\ge \rho_0),
\]
并且
\[
p'(\rho_0)=48+24\sqrt2>0,
\qquad
p(\rho_0)=16+16\sqrt2>0.
\]
因此 \(p'(\rho)>0\) 且 \(p(\rho)>0\) 对所有 \(\rho>\rho_0\) 成立，故 \(\alpha'(\rho)>0\)。
再由
\[
\alpha(\rho_0)=0,
\qquad
\lim_{\rho\to+\infty}\alpha(\rho)=\frac12,
\]
可知 \(\alpha\) 把 \((\rho_0,\infty)\) 严格单调地双射到 \((0,\tfrac12)\)。
两条单调双射合并后，唯一性与全局坐标化结论立即成立。上述导数因子分解与端点值可由脚本 \texttt{scripts/exp\_real\_input\_40\_collision\_frequency\_prime\_shadow\_audit.py} 逐项复算。证毕。
\end{proof}

\begin{corollary}[RH/Ramanujan 窗口的碰撞频率端点为五次代数数]\label{cor:xi-time-part9p-rh-window-frequency-endpoint}
令 \(u_R>1\) 为命题 \ref{prop:real-input-40-collision-rhk-window} 给出的唯一临界参数，并记
\[
\alpha_R:=\alpha(\log u_R).
\]
则 \(\alpha_R\) 是区间 \((0,\tfrac12)\) 内唯一实根，满足五次方程
\begin{equation}\label{eq:xi-time-part9p-alphaR-quintic}
191\alpha^5-701\alpha^4+1043\alpha^3-748\alpha^2+252\alpha-28=0.
\end{equation}
数值上，
\[
\alpha_R\approx 0.204230263243173.
\]
并且 RH/Ramanujan 窗口中的全部可实现碰撞频率恰为
\[
[0,\alpha_R].
\]
\end{corollary}
\begin{proof}
命题 \ref{prop:real-input-40-collision-rhk-window} 已给出 \(u_R\) 是核版 RH/Ramanujan 条件在 \(u>0\) 上的唯一正端点。另一方面，由定理
\ref{thm:xi-time-part9p-positive-branch-global-coordinate}，\(\alpha\) 作为 \(u>0\) 的函数严格递增，并把 \((0,\infty)\) 双射到 \((0,\tfrac12)\)；因此 \(u\in(0,u_R]\) 所对应的频率闭包恰为 \([0,\alpha_R]\)。

为求 \(\alpha_R\) 的代数方程，把
\eqref{eq:xi-time-part9p-collision-frequency-algebraic-curve} 与五次方程
\[
u^5+4u^4+3u^3-96u^2-42u-14=0
\]
对 \(u\) 消元。所得 \(15\) 次多项式分解为三个五次因子；其中与正实 Perron 分支兼容、并且在 \((0,\tfrac12)\) 内有实根者，恰为
\eqref{eq:xi-time-part9p-alphaR-quintic}。由上一段的单调性，该根只能是 \(\alpha_R\)，并且在 \((0,\tfrac12)\) 内唯一。数值近似由同一代数方程直接求根即得。证毕。
\end{proof}

\begin{theorem}[奇素数扭曲的 prime-shadow Euler 乘积绝对收敛]\label{thm:xi-time-part9p-odd-prime-shadow-product}
沿用命题 \ref{prop:real-input-40-collision-pressure} 的记号，并记
\[
\lambda:=\rho(1)=\varphi^2,
\qquad
t_p:=\frac{2\pi}{p}\qquad (p\ge 3\ \text{为奇素数}).
\]
由于 \(t_p<2\pi/3<R_\theta\)（推论 \ref{cor:real-input-40-collision-branch-radius}），Perron 分支 \(\rho(e^{it_p})\) 沿实轴良定义。定义
\[
\mathfrak P_{\mathrm{odd}}
:=
\prod_{p\ge 3}
\left|
\frac{\rho(e^{it_p})}{\lambda}
\right|.
\]
则该乘积绝对收敛到一个正实常数，并且
\begin{equation}\label{eq:xi-time-part9p-odd-prime-shadow-expansion}
\log \mathfrak P_{\mathrm{odd}}
=
-\frac{6\sqrt5-5}{250}(2\pi)^2P_{\mathrm{odd}}(2)
+\frac{7+24\sqrt5}{75000}(2\pi)^4P_{\mathrm{odd}}(4)
+O\!\bigl(P_{\mathrm{odd}}(6)\bigr),
\end{equation}
其中
\[
P_{\mathrm{odd}}(s):=\sum_{p\ge 3}p^{-s}.
\]
特别地，
\[
0<\mathfrak P_{\mathrm{odd}}<1.
\]
\end{theorem}
\begin{proof}
由命题 \ref{prop:arity-collision-quadratic-closed}，当 \(t\to 0\) 时有
\[
\log\left|\frac{\rho(e^{it})}{\lambda}\right|
=
-\frac{6\sqrt5-5}{250}\,t^2
+\frac{7+24\sqrt5}{75000}\,t^4
+O(t^6).
\]
取 \(t=t_p=2\pi/p\) 即得
\[
\log\left|\frac{\rho(e^{it_p})}{\lambda}\right|
=
-\frac{6\sqrt5-5}{250}\left(\frac{2\pi}{p}\right)^2
+\frac{7+24\sqrt5}{75000}\left(\frac{2\pi}{p}\right)^4
+O(p^{-6}).
\]
由于 \(\sum_{p\ge 3}p^{-2}\)、\(\sum_{p\ge 3}p^{-4}\) 与 \(\sum_{p\ge 3}p^{-6}\) 全部收敛，级数
\[
\sum_{p\ge 3}\log\left|\frac{\rho(e^{it_p})}{\lambda}\right|
\]
绝对收敛，并且逐项求和即得到
\eqref{eq:xi-time-part9p-odd-prime-shadow-expansion}。因此 Euler 乘积收敛到某个正实常数。

为证明严格上界 \(\mathfrak P_{\mathrm{odd}}<1\)，注意 \(|W(e^{it})|=A\) 且 \(A\) 为原始非负矩阵。由谱半径比较，矩阵 \(W(e^{it})\) 的谱半径不超过 \(\rho(A)=\lambda\)，故作为其本征值的 \(\rho(e^{it})\) 满足
\[
\left|\rho(e^{it})\right|\le \lambda.
\]
若某个 \(t\notin 2\pi\ZZ\) 取等，则 Wielandt 型等号情形迫使 \(W(e^{it})\) 与 \(A\) 只差一个全局相位及对角酉共轭。比较长度 \(1\) 的闭路 \(00\to 00\) 与长度 \(2\) 的闭路 \(00\to 11\to 00\) 可知该全局相位必为 \(1\)，同时又必须满足 \(e^{it}=1\)，矛盾。故对全部奇素数 \(p\) 都有
\[
\left|\rho(e^{it_p})\right|<\lambda.
\]
于是每个乘子都位于 \((0,1)\)，从而极限常数亦满足 \(0<\mathfrak P_{\mathrm{odd}}<1\)。证毕。
\end{proof}

\begin{theorem}[Rényi--2 几何下的均匀化障碍]\label{thm:xi-time-part9p-renyi2-obstruction}
\leanverified{paper\_xi\_time\_part9p\_renyi2\_obstruction}
令 \(p_m^{\mathrm{bin}}\) 为 bin-fold 推前分布，\(u_m\) 为 \(X_m\) 上的均匀分布。则 Rényi--\(2\) 散度满足精确恒等式
\[
D_2\bigl(p_m^{\mathrm{bin}}\Vert u_m\bigr)
=
\log\!\bigl(F_{m+2}\,\mathsf{Col}_m\bigr),
\]
并且
\[
D_2\bigl(p_m^{\mathrm{bin}}\Vert u_m\bigr)
\ge
\log\!\bigl(1+2C_\varphi^2\bigr)+o(1).
\]
因此
\[
\liminf_{m\to\infty}D_2\bigl(p_m^{\mathrm{bin}}\Vert u_m\bigr)
\ge
\log\!\bigl(1+2C_\varphi^2\bigr)
>
0,
\]
从而 bin-fold 推前分布在 \(\chi^2\)/Rényi--\(2\) 意义下不可能趋于均匀分布。
\end{theorem}
\begin{proof}
由 \(u_m(w)\equiv 1/F_{m+2}\) 与 Rényi--\(2\) 散度定义，
\[
D_2\bigl(p_m^{\mathrm{bin}}\Vert u_m\bigr)
=
\log\sum_{w\in X_m}\frac{\bigl(p_m^{\mathrm{bin}}(w)\bigr)^2}{u_m(w)}
=
\log\!\left(F_{m+2}\sum_{w\in X_m}\bigl(p_m^{\mathrm{bin}}(w)\bigr)^2\right).
\]
而
\[
\sum_{w\in X_m}\bigl(p_m^{\mathrm{bin}}(w)\bigr)^2
=
\frac{1}{2^{2m}}\sum_{w\in X_m}\bigl(d_m^{\mathrm{bin}}(w)\bigr)^2
=
\mathsf{Col}_m,
\]
故第一条恒等式成立。

再由推论 \ref{cor:fold-critical-resonance-gap}，
\[
\mathsf{Col}_m
\ge
\frac{1+2C_\varphi^2+o(1)}{F_{m+2}}.
\]
代回即可得到
\[
D_2\bigl(p_m^{\mathrm{bin}}\Vert u_m\bigr)
\ge
\log\!\bigl(1+2C_\varphi^2+o(1)\bigr)
=
\log\!\bigl(1+2C_\varphi^2\bigr)+o(1).
\]
取下极限便得结论。证毕。
\end{proof}

\begin{proposition}[Fibonacci 共振双频给出二阶碰撞的最小刚性证书]\label{prop:xi-time-part9p-fibonacci-resonance-cut-collision}
沿用定义 \ref{def:fold-normalized-amplitude} 的记号，
\[
a_m(k):=\frac{|\widehat d_m(k)|}{2^m},
\qquad
\mathsf{Col}_m=\frac{1}{F_{m+2}}\sum_{k\in\ZZ/(F_{m+2}\ZZ)}a_m(k)^2.
\]
定义去掉两条 Fibonacci 共振频率后的截断碰撞量
\[
\mathsf{Col}_m^{\mathrm{cut}}
:=
\frac{1}{F_{m+2}}
\sum_{\substack{k\in\ZZ/(F_{m+2}\ZZ)\\ k\notin\{F_m,F_{m+1}\}}}
a_m(k)^2.
\]
则有精确差值律
\[
\mathsf{Col}_m-\mathsf{Col}_m^{\mathrm{cut}}
=
\frac{a_m(F_m)^2+a_m(F_{m+1})^2}{F_{m+2}}
=
\frac{2C_\varphi^2+o(1)}{F_{m+2}}.
\]
因此，任何忽略 \(\{F_m,F_{m+1}\}\) 的二阶碰撞估计都会在主尺度上损失一个不可恢复的常数份额 \(2C_\varphi^2\)。
\end{proposition}
\begin{proof}
由 \(\mathsf{Col}_m\) 与 \(\mathsf{Col}_m^{\mathrm{cut}}\) 的定义，二者之差恰为被删去的两项：
\[
\mathsf{Col}_m-\mathsf{Col}_m^{\mathrm{cut}}
=
\frac{a_m(F_m)^2+a_m(F_{m+1})^2}{F_{m+2}}.
\]
再由定理 \ref{thm:fold-critical-resonance-constant}，
\[
a_m(F_m)\to C_\varphi,
\qquad
a_m(F_{m+1})\to C_\varphi.
\]
将其代回即得
\[
\mathsf{Col}_m-\mathsf{Col}_m^{\mathrm{cut}}
=
\frac{2C_\varphi^2+o(1)}{F_{m+2}}.
\]
证毕。
\end{proof}

\begin{remark}[碰撞频率端点与奇素数 shadow 的数值锚点]\label{rem:xi-time-part9p-frequency-prime-shadow-audit}
对 \(P\ge 3\)，记截断 odd-prime zeta 和
\[
P_{\mathrm{odd}}^{(\le P)}(s):=\sum_{\substack{3\le p\le P\\ p\ \mathrm{prime}}}p^{-s}.
\]
脚本 \texttt{scripts/exp\_real\_input\_40\_collision\_frequency\_prime\_shadow\_audit.py} 对 \(\alpha_R\) 与 odd-prime partial product 给出如下核验片段：
% Auto-generated; do not edit by hand.
\begin{equation}\label{eq:real_input_40_collision_frequency_prime_shadow_audit}
\begin{aligned}
\alpha_R &\approx 0.204230263243172496,\\
191\alpha_R^5-701\alpha_R^4+1043\alpha_R^3-748\alpha_R^2+252\alpha_R-28 &\approx 1.3494013e-79,\\
\mathfrak P_{\mathrm{odd}}^{(\le 2000)}:=\prod_{\substack{3\le p\le 2000\\ p\ \mathrm{prime}}}\left|\frac{\rho(e^{2\pi i/p})}{\varphi^2}\right| &\approx 0.787090286807732362,\\
\log \mathfrak P_{\mathrm{odd}}^{(\le 2000)} &\approx -0.239412314390460902,\\
-\frac{6\sqrt5-5}{250}(2\pi)^2P_{\mathrm{odd}}^{(\le 2000)}(2)+\frac{7+24\sqrt5}{75000}(2\pi)^4P_{\mathrm{odd}}^{(\le 2000)}(4) &\approx -0.250451591524040363.
\end{aligned}
\end{equation}

\end{remark}

\paragraph{Fourier conductor 壳层、真商模下降与 full-modulus 原始性}
在折叠多重度剖面的 Parseval 恒等式、群环乘积分解以及 Fibonacci 共振双频的常数级下界已经建立之后，还可把“是否能够下降到真商模”完全转写为 Fourier 支撑的 conductor 截断问题。

\begin{theorem}[折叠多重度剖面下降到真商模的 Fourier--conductor 充要条件]\label{thm:xi-fold-multiplicity-quotient-descent-fourier-conductor}
记
\[
M:=F_{m+2},
\qquad
d_m:\ZZ/(M\ZZ)\to\NN,
\qquad
\widehat d_m(k):=\sum_{r\in\ZZ/(M\ZZ)}d_m(r)e^{-2\pi i kr/M}.
\]
固定约数 \(D\mid M\)。则下列三条命题等价：
\begin{enumerate}
  \item 存在函数
  \[
  f_D:\ZZ/(D\ZZ)\to\NN
  \]
  使得
  \[
  d_m(r)=f_D(r\bmod D)
  \qquad
  (\forall r\in\ZZ/(M\ZZ));
  \]
  \item 对每个频率 \(k\in\ZZ/(M\ZZ)\)，只要
  \[
  k\not\equiv 0\pmod{M/D},
  \]
  就有
  \[
  \widehat d_m(k)=0;
  \]
  \item 若定义导出模
  \[
  \operatorname{cond}(k):=\frac{M}{\gcd(k,M)}
  \]
  及各壳层能量
  \[
  E_m(q):=\sum_{\operatorname{cond}(k)=q}|\widehat d_m(k)|^2
  \qquad (q\mid M),
  \]
  则
  \[
  E_m(q)=0
  \qquad
  \text{对一切 }q\nmid D.
  \]
\end{enumerate}
此外，不加任何额外假设时总有精确分解
\[
\mathsf{Col}_m
=
\frac1M+\frac{1}{2^{2m}M}\sum_{\substack{q\mid M\\ q>1}}E_m(q),
\]
而在上述等价条件成立时，上式退化为
\[
\mathsf{Col}_m
=
\frac1M+\frac{1}{2^{2m}M}\sum_{\substack{q\mid D\\ q>1}}E_m(q).
\]
\end{theorem}
\begin{proof}
先证 \((1)\Rightarrow(2)\)。若
\[
d_m(r)=f_D(r\bmod D),
\]
则把 \(r\) 唯一写成
\[
r=a+Db,
\qquad
0\le a<D,
\qquad
0\le b<\frac{M}{D}.
\]
于是
\[
\begin{aligned}
\widehat d_m(k)
&=
\sum_{a=0}^{D-1}\sum_{b=0}^{M/D-1}
f_D(a)\,e^{-2\pi i k(a+Db)/M}\\
&=
\left(\sum_{a=0}^{D-1}f_D(a)e^{-2\pi i ka/M}\right)
\left(\sum_{b=0}^{M/D-1}e^{-2\pi i kb/(M/D)}\right).
\end{aligned}
\]
第二个括号是长度 \(M/D\) 的几何和；当且仅当
\[
\frac{M}{D}\mid k
\]
时非零，否则为零。故 \((2)\) 成立。

再证 \((2)\Rightarrow(1)\)。Fourier 反演给出
\[
d_m(r)=\frac1M\sum_{k\in\ZZ/(M\ZZ)}\widehat d_m(k)e^{2\pi i kr/M}.
\]
在 \((2)\) 的假设下，只保留 \(k=(M/D)\ell\) 的项，可得
\[
\begin{aligned}
d_m(r)
&=
\frac1M\sum_{\ell=0}^{D-1}\widehat d_m\!\left(\frac{M}{D}\ell\right)
e^{2\pi i \ell r/D}.
\end{aligned}
\]
右端仅依赖于 \(r\bmod D\)，故 \(d_m\) 下降到 \(\ZZ/(D\ZZ)\)。

现证 \((2)\Leftrightarrow(3)\)。若
\[
k=\frac{M}{D}\ell,
\]
则
\[
\gcd(k,M)=\frac{M}{D}\gcd(\ell,D),
\qquad
\operatorname{cond}(k)=\frac{D}{\gcd(\ell,D)}\mid D.
\]
反之，若 \(\operatorname{cond}(k)\mid D\)，记
\[
g:=\gcd(k,M),
\qquad
\operatorname{cond}(k)=\frac{M}{g}.
\]
由 \(\frac{M}{g}\mid D\) 且 \(D\mid M\)，可写
\[
D=\frac{M}{g}\,t
\]
对某个整数 \(t\)。于是
\[
\frac{M}{D}=\frac{g}{t},
\]
从而 \(\frac{M}{D}\mid g\)。由于 \(g\mid k\)，遂有
\[
\frac{M}{D}\mid k.
\]
故
\[
\frac{M}{D}\mid k
\iff
\operatorname{cond}(k)\mid D,
\]
于是 \((2)\) 与 \((3)\) 等价。

最后证明碰撞分解。由定理 \ref{thm:fold-collision-spectrum}，
\[
\mathsf{Col}_m
=
\frac{1}{2^{2m}M}\sum_{k\in\ZZ/(M\ZZ)}|\widehat d_m(k)|^2.
\]
其中 \(k=0\) 一项恰给出 \(1/M\)，而非零频率按 \(\operatorname{cond}(k)\) 分组，便得到
\[
\mathsf{Col}_m
=
\frac1M+\frac{1}{2^{2m}M}\sum_{\substack{q\mid M\\ q>1}}E_m(q).
\]
若 \((1)\)--\((3)\) 成立，则所有 \(q\nmid D\) 的壳层均为零，故立即退化为
\[
\mathsf{Col}_m
=
\frac1M+\frac{1}{2^{2m}M}\sum_{\substack{q\mid D\\ q>1}}E_m(q).
\]
证毕。
\end{proof}
\leanverified{paper\_xi\_fold\_multiplicity\_quotient\_descent\_fourier\_conductor}

\begin{corollary}[大尺度下折叠多重度剖面的 full-modulus 原始性]\label{cor:xi-fold-multiplicity-fullmodulus-primitive}
存在整数 \(m_0\)，使得对所有 \(m\ge m_0\)，折叠多重度剖面
\[
d_m:\ZZ/(F_{m+2}\ZZ)\to\NN
\]
都不可能下降到任何真商模
\[
\ZZ/(D\ZZ),
\qquad
D\mid F_{m+2},
\qquad
D<F_{m+2}.
\]
等价地，对充分大的 \(m\)，其最小可见模数正是 \(F_{m+2}\) 本身。
\end{corollary}
\begin{proof}
由定理 \ref{thm:fold-critical-resonance-constant}，
\[
a_m(F_m):=\frac{|\widehat d_m(F_m)|}{2^m}\longrightarrow C_\varphi>0.
\]
故存在 \(m_0\)，使得当 \(m\ge m_0\) 时，
\[
\widehat d_m(F_m)\neq 0.
\]
另一方面，由 Fibonacci 最大公因子恒等式，
\[
\gcd(F_m,F_{m+2})=F_{\gcd(m,m+2)}=1,
\]
从而
\[
\operatorname{cond}(F_m)=\frac{F_{m+2}}{\gcd(F_m,F_{m+2})}=F_{m+2}.
\]
若 \(d_m\) 能下降到某个真商模 \(D<F_{m+2}\)，则由定理
\ref{thm:xi-fold-multiplicity-quotient-descent-fourier-conductor} 的条件 \((3)\)，所有
\[
q\nmid D
\]
的 conductor 壳层都必须消失。特别地，full-conductor 壳层
\[
q=F_{m+2}
\]
必须为零，这与 \(\widehat d_m(F_m)\neq 0\) 矛盾。故不存在这样的真商模 \(D\)。证毕。
\end{proof}

\begin{theorem}[full-conductor 壳层的碰撞缺口下界由 Fibonacci 共振双频承担]\label{thm:xi-full-conductor-shell-collision-gap-fibonacci-pair}
记
\[
M:=F_{m+2},
\qquad
E_m^{\mathrm{full}}:=E_m(M)=\sum_{\gcd(k,M)=1}|\widehat d_m(k)|^2,
\qquad
a_m(k):=\frac{|\widehat d_m(k)|}{2^m}.
\]
则有
\[
E_m^{\mathrm{full}}
\ge
|\widehat d_m(F_m)|^2+|\widehat d_m(F_{m+1})|^2
=
2^{2m}\bigl(a_m(F_m)^2+a_m(F_{m+1})^2\bigr).
\]
从而
\[
\mathsf{Col}_m-\frac1M
\ge
\frac{E_m^{\mathrm{full}}}{2^{2m}M}
\ge
\frac{a_m(F_m)^2+a_m(F_{m+1})^2}{M}
=
\frac{2C_\varphi^2+o(1)}{M}.
\]
特别地，碰撞缺口的主尺度下界已经完全由 full-conductor 壳层内的共轭 Fibonacci 双频承担。
\end{theorem}
\begin{proof}
由定理 \ref{thm:xi-fold-multiplicity-quotient-descent-fourier-conductor} 的壳层分解，
\[
\mathsf{Col}_m-\frac1M
=
\frac{1}{2^{2m}M}\sum_{\substack{q\mid M\\ q>1}}E_m(q)
\ge
\frac{E_m^{\mathrm{full}}}{2^{2m}M}.
\]

又由 Fibonacci 最大公因子恒等式，
\[
\gcd(F_m,F_{m+2})=\gcd(F_{m+1},F_{m+2})=1,
\]
故频率 \(F_m\) 与 \(F_{m+1}\) 都落在 full-conductor 壳层内。于是
\[
E_m^{\mathrm{full}}
\ge
|\widehat d_m(F_m)|^2+|\widehat d_m(F_{m+1})|^2.
\]
再用
\[
|\widehat d_m(F_m)|=2^m a_m(F_m),
\qquad
|\widehat d_m(F_{m+1})|=2^m a_m(F_{m+1}),
\]
便得到第二个不等式。

最后，由定理 \ref{thm:fold-critical-resonance-constant}，
\[
a_m(F_m)\to C_\varphi,
\qquad
a_m(F_{m+1})\to C_\varphi,
\]
故
\[
\frac{a_m(F_m)^2+a_m(F_{m+1})^2}{M}
=
\frac{2C_\varphi^2+o(1)}{M}.
\]
证毕。
\end{proof}

\begin{theorem}[Fourier 幅度的中点门级联实现与临界共振半径]\label{thm:xi-time-part9p1-midpoint-gate-cascade-critical-radius}
沿用定理 \ref{thm:fold-spectrum-factorization} 的记号。对每个 \(m\ge 1\) 与 \(k\in\ZZ/(F_{m+2}\ZZ)\)，定义角序列
\[
\vartheta_j^{(m)}(k):=
2\pi k\frac{F_{j+1}}{F_{m+2}}
\qquad (1\le j\le m),
\]
以及复平面上的中点门
\[
\mathcal M_{\vartheta}(z):=\frac{z+e^{-\mathrm{i}\vartheta}z}{2}.
\]
再置
\[
z_0^{(m,k)}:=1,
\qquad
z_j^{(m,k)}:=\mathcal M_{\vartheta_j^{(m)}(k)}\bigl(z_{j-1}^{(m,k)}\bigr)
\qquad (1\le j\le m).
\]
则成立：
\begin{enumerate}
  \item 级联终值满足
  \[
  \boxed{\;
  z_m^{(m,k)}
  =
  2^{-m}\widehat d_m(k),\;}
  \]
  从而
  \[
  \boxed{\;
  |z_m^{(m,k)}|=a_m(k).\;}
  \]
  因此 \(a_m(k)\) 精确等同于角序列 \(\{\vartheta_j^{(m)}(k)\}_{j=1}^{m}\) 所驱动的中点门级联收缩半径；
  \item 有
  \[
  \widehat d_m(k)=0
  \iff
  z_m^{(m,k)}=0
  \iff
  \exists\, j\in\{1,\dots,m\}\ \text{使}\ 
  \vartheta_j^{(m)}(k)\equiv \pi\pmod{2\pi}.
  \]
  换言之，Fourier 零点恰对应于某一级中点门把一点与其对径旋转副本取平均而坍缩到原点；
  \item 对 Fibonacci 补频对 \(k=F_m\) 与 \(k=F_{m+1}\)，有
  \[
  |z_m^{(m,F_m)}|\longrightarrow C_\varphi,
  \qquad
  |z_m^{(m,F_{m+1})}|\longrightarrow C_\varphi.
  \]
  因而临界共振常数 \(C_\varphi\) 可解释为一条中点门级联在 Fibonacci 角驱动下保持的严格正极限半径。
\end{enumerate}
\end{theorem}
\begin{proof}
由定义，
\[
\mathcal M_{\vartheta}(z)=\frac{1+e^{-\mathrm{i}\vartheta}}{2}\,z.
\]
故递推展开给出
\[
z_m^{(m,k)}
=
\prod_{j=1}^{m}\frac{1+e^{-\mathrm{i}\vartheta_j^{(m)}(k)}}{2}
=
2^{-m}\prod_{j=1}^{m}\left(1+e^{-2\pi \mathrm{i}kF_{j+1}/F_{m+2}}\right).
\]
再由定理 \ref{thm:fold-spectrum-factorization}，
\[
\widehat d_m(k)=\prod_{j=1}^{m}\left(1+e^{-2\pi \mathrm{i}kF_{j+1}/F_{m+2}}\right),
\]
便得到
\[
z_m^{(m,k)}=2^{-m}\widehat d_m(k).
\]
取模即得
\[
|z_m^{(m,k)}|=\frac{|\widehat d_m(k)|}{2^m}=a_m(k).
\]

第二条由乘积分解立即推出：\(\widehat d_m(k)=0\) 当且仅当某一因子
\[
1+e^{-\mathrm{i}\vartheta_j^{(m)}(k)}
\]
为零，而这又等价于
\[
\vartheta_j^{(m)}(k)\equiv \pi\pmod{2\pi}.
\]
此时
\[
\mathcal M_{\vartheta_j^{(m)}(k)}(z)=0
\qquad (\forall z\in\CC),
\]
故级联在该步后永久坍缩到原点。

第三条正是定理 \ref{thm:fold-critical-resonance-constant} 的结论，因为
\[
|z_m^{(m,F_m)}|=a_m(F_m),
\qquad
|z_m^{(m,F_{m+1})}|=a_m(F_{m+1}).
\]
证毕。
\end{proof}

\endinput

